% =====================================================================
%  Chapter 31 station file -- 14 chaptersections + Summit + lastword.
%    1 Summary                8 Recommended Study Order
%    2 Master Decision Flow   9 Difficulty Progression
%    3 Common Mistakes       10 Reduction Gallery
%    4 WA Clinic             11 Mixed Problem Set (10 practice items)
%    5 Quick Reference       12 Summit -- CF 888G Xor-MST
%    6 Write-From-Memory     13 Tool Kit Inventory
%    7 Complete Checklist    14 Related Topics
%  Summit live-verified: title "G. Xor-MST", 2s TL, 256 MB, n<=200000,
%  a_i < 2^30.  Two official samples: [1 2 3 4 5] -> 8; [1 2 3 4] -> 8.
%  Both hand-replayed in the Trace section via the bit divide-and-conquer.
%  Mixed-set IDs live-verified (title + rating) and absent from used_cf.txt;
%  \practice uses the correct FIVE-argument form (title / judge / url /
%  difficulty / hint), matching Ch.29, NOT Ch.30's leaked seven-arg form.
% =====================================================================

\chaptersection{Chapter Summary}

Chapter~31 is the optimisation counterpart of Chapter~30's enumeration.  Where
\chref{c32} counted spanning trees with a determinant, this chapter picks the
\emph{cheapest} one and interrogates its structure.  Four tools carry it:

\begin{itemize}[nosep]
  \item the \textbf{cut property} (\(\thmref{c17-cut}\)): the lightest edge across
        any cut lies in some MST -- the single lemma behind every greedy here;
  \item \textbf{Kruskal} (\(\secref{c17-kruskal}\)) and \textbf{Prim}
        (\(\secref{c17-prim}\)): sort-and-union with a DSU (\(\bigO{m\log m}\))
        versus grow-one-frontier with a heap (\(\bigO{m\log n}\));
  \item the \textbf{Kruskal-reconstruction tree} (\(\secref{c17-krt}\)): record
        each union as a node, and a bottleneck path query becomes an LCA weight
        (\(\thmref{c17-krtmax}\), \(\chref{c11}\));
  \item the \textbf{edge classification} (\(\secref{c17-classify}\)): forced /
        optional / forbidden via the intersection (\(\thmref{c17-intersect}\))
        and cycle (\(\thmref{c17-cycle}\)) properties, counted per equal-weight
        group (\(\lemref{c17-group}\)).
\end{itemize}

\(\secref{c17-boundary}\) adds the second-best MST (\(\thmref{c17-second}\)) by a
single replacement swap and draws the hand-off line: counting is Ch.~30,
unweighted robustness is Ch.~29, directed arborescences are out of scope.  The
Summit \(\probref{c17-summit}\) (CF 888G) is a complete graph whose edge weight is
an xor -- the cut property applied over value bits by a trie divide-and-conquer.

\chaptersection{Master Decision Flowchart}

\begin{center}
\begin{tikzpicture}[diagram, scale=0.82, transform shape,
   q/.style={rectangle, rounded corners=3pt, draw=primaryblue!70, fill=primaryblue!8,
             line width=0.8pt, minimum width=24mm, minimum height=8mm,
             font=\scriptsize\bfseries, inner sep=2pt, align=center},
   a/.style={rectangle, rounded corners=2pt, draw=successgreen!70, fill=successgreen!10,
             line width=0.8pt, minimum width=30mm, minimum height=8mm,
             font=\scriptsize, inner sep=2pt, align=center},
   every edge/.style={-Latex, line width=0.8pt, draw=darkgray}]
  \node[q] (root) at (0,0) {weighted, undirected;\\ what is asked?};
  \node[q] (w)  at (-5.5,-1.7) {MST weight\\ or the tree};
  \node[q] (q)   at (0,-1.7) {path / bottleneck\\ questions};
  \node[q] (s)   at (5.5,-1.7) {which edges\\ / uniqueness};
  \node[a] (kru) at (-7.5,-3.4) {sparse: Kruskal\\ \secref{c17-kruskal}};
  \node[a] (pri) at (-4,-3.4) {dense: Prim\\ \secref{c17-prim}};
  \node[a] (krt) at (-0.2,-3.4) {reconstruction\\ tree \secref{c17-krt}};
  \node[a] (cla) at (3.6,-3.4) {group count\\ \secref{c17-classify}};
  \node[a] (sec) at (7.4,-3.4) {second best\\ \secref{c17-boundary}};
  \node[a] (bo)  at (0,-5) {count -> Ch.30; bridges -> Ch.29; directed -> out};
  \draw (root)--(w)--(kru);
  \draw (w)--(pri);
  \draw (root)--(q)--(krt);
  \draw (root)--(s)--(cla);
  \draw (s)--(sec);
  \draw (q.south) to[out=-90,in=90] (bo.north);
\end{tikzpicture}
\end{center}

Every Ch.~31 problem walks one of the six green boxes.  The bottom box is
\(\secref{c17-boundary}\)'s pick list: a question that says "cheapest edges" but
means a count (Ch.~30), a bridge (Ch.~29), or a directed arborescence is not an
MST problem.

\chaptersection{Common Mistakes Checklist}

\begin{enumerate}[nosep]
  \item \textbf{"Locally cheapest is globally safe."}  Only an edge that is the
        cheapest across the cut defined by the \emph{current} components is safe
        (\(\thmref{c17-cut}\)); a raw greedy over unsorted edges is wrong.
  \item \textbf{Prim keyed by accumulated distance.}  Prim's frontier key is the
        single edge weight, not a source distance -- that is Dijkstra
        (\(\secref{c17-prim}\)).
  \item \textbf{Sorting only once for a dynamic MST.}  Kruskal's sort is a one-shot;
        re-sorting per query is a \(\bigO{m\log m}\) tax you can avoid with a
        reconstruction tree.
  \item \textbf{Forgetting the pre-union count.}  In the equal-weight group loop
        (\(\lemref{c17-group}\)), \(\mathit{usable}\) must be read \emph{before}
        any union in the block; a single-scan merge undercounts and fails CF 1108F.
  \item \textbf{No parent / rank in the DSU.}  Kruskal leans on \(\chref{c28}\);
        a path-compression-less DSU on \(m=2\cdot10^5\) can degrade.
  \item \textbf{Int overflow on MST sums.}  CF 1095F edges reach \(2\cdot10^{12}\);
        accumulate in \texttt{long long} (\(\secref{c17-kruskal}\)).
  \item \textbf{Treating a second-best swap as two swaps.}  \(\thmref{c17-second}\)
        needs only \emph{one} non-tree edge added and one tree edge dropped.
  \item \textbf{Feeding a directed graph to Kruskal.}  Directed MST is Edmonds,
        not this chapter (\(\secref{c17-boundary}\)).
\end{enumerate}

\chaptersection{Wrong-Answer Clinic}

\begin{center}
\begin{tikzpicture}[diagram, scale=0.82, transform shape,
   r/.style={rectangle, rounded corners=3pt, draw=warningred!70, fill=warningred!8,
             line width=0.9pt, minimum width=22mm, minimum height=8mm,
             font=\scriptsize\bfseries, inner sep=2pt, align=center},
   f/.style={rectangle, rounded corners=2pt, draw=successgreen!70, fill=successgreen!10,
             line width=0.8pt, minimum width=40mm, minimum height=9mm,
             font=\scriptsize, inner sep=2pt, align=center},
   every edge/.style={-Latex, line width=0.8pt, draw=darkgray}]
  \node[r] (v) at (0,0) {verdict on the judge};
  \node[r] (wa) at (-5,-1.5) {WA};
  \node[r] (tle) at (0,-1.5) {TLE};
  \node[r] (re) at (5,-1.5) {RE / overflow};
  \node[f] (w1) at (-6.5,-3.4) {single-scan group: read usable BEFORE unions (\(\secref{c17-classify}\))};
  \node[f] (w2) at (-2.8,-3.4) {wrong anchor city in CF 1095F: pick \(\arg\min a_i\) first};
  \node[f] (t1) at (0.6,-3.4) {materialised \(\binom n2\) edges: reduce to star + offers};
  \node[f] (r1) at (4.3,-3.4) {int sum overflow: use \texttt{long long}};
  \node[f] (r2) at (7,-3.4) {deep DSU: add path compression};
  \draw (v)--(wa)--(w1);  \draw (wa)--(w2);
  \draw (v)--(tle)--(t1);
  \draw (v)--(re)--(r1);  \draw (re)--(r2);
\end{tikzpicture}
\end{center}

\begin{itemize}[nosep]
  \item \textbf{WA on CF 1108F.}  A single-scan group loop: the second edge of a
        same-weight block sees an already-merged component and is not counted
        usable, so the excess is under-tallied.  Split into a count-then-merge
        pair of loops (\(\coderef{c17-group}\)).
  \item \textbf{WA on CF 1095F.}  Anchoring the star at the wrong city (not the
        minimum \(a_i\)) inflates every attachment; the dominance argument needs
        \(s=\arg\min a_i\) (\(\probref{c17-connect}\)).
  \item \textbf{TLE on CF 1095F or 888G.}  Building the implicit complete graph
        (\(\binom n2\) edges) instead of the \(O(n)\) candidate set -- the star
        plus offers, or the trie recursion.
  \item \textbf{Overflow on CF 1095F.}  Weights up to \(2\cdot10^{12}\) and totals
        to \(\sim2\cdot10^{17}\): an \texttt{int} accumulator wraps to a wrong answer.
\end{itemize}

\chaptersection{Quick Reference Card}

\begin{quickref}[Chapter 31 at a glance]
\begin{itemize}[nosep]
  \item \textbf{Cut property.}  Lightest edge across any cut is in some MST; if
        strictly lightest, in every MST (\(\thmref{c17-cut}\)).
  \item \textbf{Kruskal.}  Sort \(+\) DSU, \(\bigO{m\log m}\).  Prim (heap)
        \(\bigO{m\log n}\); Prim (matrix) \(\bigO{n^2}\).
  \item \textbf{Reconstruction tree.}  Union node weight \(=\) edge; bottleneck
        \(\mathrm{maxw}(u,v)=\mathrm{val}(\lca(u,v))\) (\(\thmref{c17-krtmax}\)).
  \item \textbf{Classification.}  forced / optional / forbidden via
        \(\thmref{c17-intersect}\), \(\thmref{c17-cycle}\).
  \item \textbf{Group excess.}  \(\sum_{\text{blocks}}(\mathit{usable}-\mathit{merged})\)
        optional edges; equals CF 1108F's answer (\(\lemref{c17-group}\)).
  \item \textbf{CF 1095F.}  Star to \(\arg\min a_i\) (weights \(a_i+a_s\)) plus the
        \(m\) offers, Kruskal.
  \item \textbf{Second-best.}  \(\min_{e\notin T^{*}} w(T^{*})+w(e)-\mathrm{maxw}_{T^{*}}(e)\)
        (\(\thmref{c17-second}\)).
  \item \textbf{Summit CF 888G.}  \(\mathrm{MST}(S)=\mathrm{solve}(L)+\mathrm{solve}(R)+
        \text{minCrossXor}\), trie for the cross term.
\end{itemize}
\end{quickref}

\chaptersection{Write-From-Memory Drills}

Each drill rebuilds a listing that exists in the chapter; the blanks are the
decisive lines.

\begin{enumerate}[nosep]
  \item Kruskal: sort by weight, then take the edge only when its endpoints lie in
different components and union them -- rebuild \(\coderef{c17-kruskal}\); blank:
the cycle test ("same component yet?").
  \item Prim (heap): push \((w,v)\) neighbours, pop the cheapest unseen -- rebuild
        \(\coderef{c17-prim}\); blank: the frontier key is the single edge weight
        \(w\), not \(w+\) a source distance.
  \item Reconstruction tree: on a successful union, allocate a node and adopt both
        representatives -- rebuild \(\coderef{c17-krt}\); blank: \(\mathrm{val}[\text{new}]=w\).
  \item CF 1108F group scan: two loops per equal-weight block -- rebuild
        \(\coderef{c17-group}\); blanks: count \(\mathit{usable}\) before merging,
        add \(\mathit{usable}-\mathit{merged}\).
  \item Second-best: for each non-tree edge, add \(w(T^{*})+w(e)-\) (heaviest tree
        edge on its cycle) -- rebuild \(\coderef{c17-second}\); blank: the max-on-path
        value \(\mathrm{maxw}\) via LCA / KRT (\(\chref{c11}\)).
\end{enumerate}

\chaptersection{Complete Problem Checklist}

\begin{center}
\renewcommand{\arraystretch}{1.2}\setlength{\tabcolsep}{3pt}
\begin{tabular}{@{}p{0.14\linewidth} >{\raggedright\arraybackslash}p{0.26\linewidth} >{\raggedright\arraybackslash}p{0.22\linewidth} >{\raggedright\arraybackslash}p{0.10\linewidth} >{\raggedright\arraybackslash}p{0.14\linewidth}@{}}
\toprule
\textbf{ID} & \textbf{Title} & \textbf{Tool} & \textbf{Rating} & \textbf{Section}\\
\midrule
CF 1095F & Make It Connected & implicit star + offers, Kruskal & 1900 & 31.1\\
CF 1108F & MST Unification & equal-weight group excess & 2100 & 31.4\\
CF 888G  & Xor-MST & trie divide-and-conquer + cut property & 2300 & Summit\\
\bottomrule
\end{tabular}
\end{center}

The worked examples (not graded): Kruskal and the cut property on \(G\) (31.1),
Prim on \(G\) (31.2), the reconstruction tree and three bottleneck queries on
\(G\) (31.3), and the second-best swap on \(G\) (31.5).

\chaptersection{Recommended Study Order}

Read \(\secref{c17-kruskal}\) and hand-run Kruskal on the running graph \(G\)
until the weight \(12\) is muscle memory; the cut property is the one idea every
later section leans on.  Then \(\secref{c17-prim}\) -- skim the heap version, do
the dense \(\bigO{n^2}\) matrix Prim by hand on \(G\) so the frontier view sticks,
and internalise the "Prim is not Dijkstra" warning.  Solve CF 1095F only after you
can state why the star to the cheapest city dominates.

\(\secref{c17-krt}\) is best taken once Kruskal is automatic; build the merge tree
of \(G\) by hand and check the three queries against the MST paths.  Then
\(\secref{c17-classify}\), whose group-excess is the crux of CF 1108F -- trace
sample 6 (the answer \(2\)) before coding.  Close with \(\secref{c17-boundary}\)
and the Summit CF 888G.

\chaptersection{Difficulty Progression}

\begin{itemize}[nosep]
  \item \lv{3} State the cut and cycle properties; hand-run Kruskal and Prim on
        \(G\), both giving weight \(12\).
  \item \lv{4} Code Kruskal with a DSU and dense Prim; build the reconstruction
        tree and answer a bottleneck query as an LCA weight; solve CF 1095F.
  \item \lv{5} Solve CF 1108F (group excess \(=\) increments to a unique MST);
        compute the second-best MST of \(G\) by a single swap.
  \item \lv{5} (stretch) The Summit CF 888G: the xor divide-and-conquer, both
        official samples replayed to \(8\).
\end{itemize}

\chaptersection{Reduction Gallery}

\begin{itemize}[nosep]
  \item \emph{Cheapest network of cables with a per-city hookup tax} \(\to\) an
        implicit star edge cost \(a_x+a_y\) \(\to\) reduce to star + offers,
        Kruskal (CF 1095F).
  \item \emph{Minimise the worst link on every route} \(\to\) bottleneck path on
        the MST \(\to\) reconstruction tree + LCA (\(\thmref{c17-krtmax}\)).
  \item \emph{Connect all with the fewest ties / a unique plan} \(\to\)
        equal-weight group counting \(\to\) \(\lemref{c17-group}\) (CF 1108F).
  \item \emph{Best alternative network} \(\to\) one non-tree edge in, one tree edge
        out \(\to\) second-best swap (\(\thmref{c17-second}\)).
  \item \emph{Complete graph, edge weight an xor} \(\to\) split on the top bit, one
        cross edge by the cut property \(\to\) trie greedy (\(\probref{c17-summit}\)).
  \item \emph{Threshold "smallest \(t\) so \(u,v\) connect using edges \(\le t\)"}.
        \(\to\) it is exactly \(\mathrm{val}(\lca(u,v))\) in the reconstruction tree.
\end{itemize}

\chaptersection{Mixed Problem Set}

\begin{mixedset}
\practice{Edgy Trees}{CF 1139C}{https://codeforces.com/problemset/problem/1139/C}{CF 1500 L3}{No weights here -- keep only the live links and count whole connected pieces. This is a component census, not a cheapest-edges choice; see \chref{c28}.}
\practice{Roads not only in Berland}{CF 25D}{https://codeforces.com/problemset/problem/25/D}{CF 1900 L4}{Decide which old links to keep and which to repurpose so every region reaches a capital; process the cheapest commitments first.}
\practice{Minimal Diameter Forest}{CF 1092E}{https://codeforces.com/problemset/problem/1092/E}{CF 2000 L4}{You must join already-connected pieces with a few new links so the longest distance stays small; balance each piece's reach as you merge.}
\practice{New Roads}{CF 746G}{https://codeforces.com/problemset/problem/746/G}{CF 2100 L4}{The price of a link depends on how far the endpoint already sits from the capital; choose the join order so the total cost profile is cheapest.}
\practice{Village (Minimum)}{CF 1387B1}{https://codeforces.com/problemset/problem/1387/B1}{CF 2100 L4}{The network is already a tree and nothing is being chosen by weight -- this is a cost redistribution to a meeting point; see \chref{c10}.}
\practice{Hamiltonian Spanning Tree}{CF 618D}{https://codeforces.com/problemset/problem/618/D}{CF 2200 L4}{It must visit every vertex exactly once along links that are NOT the given tree -- a Hamiltonian-path decision, not a cheap-edges spanning tree.}
\practice{Tree}{CF 932D}{https://codeforces.com/problemset/problem/932/D}{CF 2200 L4}{Edges arrive over time with weights and you answer best-downward-path questions on the growing structure -- a distance/LCA maintenance task; see \chref{c11}.}
\practice{Complete the MST}{CF 1508C}{https://codeforces.com/problemset/problem/1508/C}{CF 2500 L5}{Some links are already fixed for you; pick the remaining ones so the whole thing connects with the cheapest total.}
\practice{MST in Modulo Graph}{CF 2081D}{https://codeforces.com/problemset/problem/2081/D}{CF 2600 L5}{The graph is implicit through residues, but only a handful of candidate links can ever be the cheap way to attach a vertex -- materialise just those.}
\practice{Best Edge Weight}{CF 827D}{https://codeforces.com/problemset/problem/827/D}{CF 2700 L5}{Ask how heavy each single link may become before the cheapest network stops using it; the merge history of Kruskal carries the answer.}
\end{mixedset}

\chaptersection{The Summit -- CF 888G Xor-MST}

\begin{notebox}[Why this is the chapter top rope]
CF 888G hides an MST inside a \emph{complete} graph on up to \(2\cdot10^5\)
vertices -- \(\Theta(n^2)\) edges, so no explicit Kruskal is possible.  The edge
weight is \(a_u \oplus a_v\), and the whole MST falls out of the cut property
applied over the \textbf{value bits}: split the numbers by their highest set bit,
recurse into each half, and pay exactly \emph{one} cross edge whose xor is the
minimum across the split (found by a binary trie).  It composes two chapter tools
(the cut property and Kruskal's sorted-by-weight view) with an earlier one -- the
bottom-up divide-and-conquer fold of Ch.~13 (\chref{c10}) -- over the bit lattice.
\end{notebox}

\solvedproblem{c17-summit}{Xor-MST}{CF 888G}{https://codeforces.com/problemset/problem/888/G}{CF 2300 L5}{cut property over value bits, binary trie greedy, divide-and-conquer fold (Ch.13)}

\begin{problemstatement}[CF 888G, verified against the live CF problem page]
You are given \(n\) integers \(a_1,\dots,a_n\) (\(2\le n\le 2\cdot 10^5\),
\(0\le a_i < 2^{30}\)).  Consider the complete graph on these \(n\) vertices where
the weight of edge \((u,v)\) is \(a_u \oplus a_v\) (bitwise xor).  Find the total
weight of the minimum spanning tree of this graph.  Time limit 2 seconds, memory
256 MB.  Samples: for \([1,2,3,4,5]\) the answer is \(8\); for \([1,2,3,4]\) the
answer is \(8\).
\end{problemstatement}

\begin{figure}[H]
\centering
\begin{tikzpicture}[diagram, scale=0.9, transform shape,
   tnode/.style={circle, draw=primaryblue!70, fill=primaryblue!8, line width=0.7pt,
                 minimum size=6mm, font=\scriptsize\bfseries, inner sep=1pt}]
  \node[tnode] (a) at (0,0) {1};
  \node[tnode] (b) at (2,1.2) {2};
  \node[tnode] (c) at (4,0) {3};
  \node[tnode] (d) at (2,-1.4) {4};
  \node[tnode] (e) at (3.4,-2.6) {5};
  \draw[darkgray!40] (a)--(b)--(c)--(d)--(e)--(a) (a)--(c) (b)--(d) (b)--(e) (c)--(e) (d)--(e);
  \node[annot, font=\scriptsize, text width=6cm, align=center] at (2,-3.8)
    {the disguise: every pair is an edge, weight = xor; \(\binom n2\) edges, no explicit Kruskal};
\end{tikzpicture}
\caption{The disguise (sample 1): a complete graph whose edge weights are xors of
the vertex values.  Drawing it directly and sorting \(\binom n2\) edges is
hopeless; the reduction turns it into one recursion over bits.}
\label{fig:c17-summit-disguise}
\end{figure}

\begin{figure}[H]
\centering
\begin{tikzpicture}[diagram, scale=0.9, transform shape,
   half/.style={rectangle, rounded corners=2pt, draw=softpurple!70, fill=softpurple!10,
                line width=0.9pt, minimum width=18mm, minimum height=9mm,
                font=\scriptsize\bfseries, inner sep=2pt, align=center}]
  \node[half] (L) at (0,0) {\(L=\{1,2,3\}\)\\ bit2 = 0};
  \node[half] (R) at (5,0) {\(R=\{4,5\}\)\\ bit2 = 1};
  \draw[difforange, line width=1.6pt] (L.east)--node[fill=white,font=\scriptsize,inner sep=1.5pt] {minCross = 4} (R.west);
  \node[annot, font=\scriptsize, align=center] at (0,-1.3) {solve(L) = 3};
  \node[annot, font=\scriptsize, align=center] at (5,-1.3) {solve(R) = 1};
  \node[annot, font=\scriptsize, text width=8cm, align=center] at (2.5,-2.6)
    {MST = solve(L) + solve(R) + 1 cross edge of xor \(4\)  \(= 3 + 1 + 4 = 8\)};
\end{tikzpicture}
\caption{The true structure after the reduction (sample 1, top bit \(2^2\)).
Vertices split into \(L=\{1,2,3\}\) (bit~2 clear) and \(R=\{4,5\}\) (bit~2 set).
The MST cost is the two halves' costs plus the single cheapest cross edge
(\(1\oplus 5 = 4\)), which the cut property forces to exactly one edge.}
\label{fig:c17-summit-reduce}
\end{figure}

\begin{attemptgate}[Attempt it yourself first -- \lv{5}, give yourself 60 minutes]
You cannot materialise \(\binom n2\) xor edges.  Look at the highest bit on which
the numbers differ: what does the cut property say about how many edges must cross
that bit-split?
\peek{Group the values by their most-significant differing bit.  Any two numbers
with different top bits have an xor with that bit set -- so the cheap edges all sit
\emph{inside} each group, and by the cut property (\(\thmref{c17-cut}\)) exactly
one edge joins the two groups.  That single joining edge is the minimum xor across
the split.}
\nudge{Recurse: \(\mathrm{MST}(S) = \mathrm{MST}(L) + \mathrm{MST}(R) +\)
(minimum \(x\oplus y\) over \(x\in L, y\in R\)), splitting on the current highest
bit.  The base case (one value, or all bits consumed) costs \(0\).  The cross term
is a classic binary-trie query.}
\fullhint{Insert one half into a bitwise trie and, for each value of the other half,
greedily walk the trie preferring the same bit (to keep the xor small), accumulating
\((1\ll b)\) whenever only the opposite child exists -- the minimum xor for that
value.  Minimise over the half.  Each level moves every value into exactly one
child and the cross scans total \(\bigO{n\cdot 30}\) amortised, giving
\(\bigO{n\cdot 30}\) overall -- the same divide-and-conquer fold shape as Ch.~13
(\chref{c10}).  Full build below.}
\end{attemptgate}

\solstep{Understanding}
The complete graph is too big to see, but its MST is tiny in structure: on edge
weights that are xors, the highest differing bit dominates.  All edges inside a
same-bit group have that bit clear in their xor (hence cheap), while every
cross-group edge has it set (hence expensive by \(2^b\)).  So the MST fully builds
each group internally and pays exactly one cross edge to glue them -- the cut
property in the value-bit world.

\solstep{Key Observations}
\begin{observation}
Split \(S\) at bit \(b\) into \(L\) (bit \(b\) clear) and \(R\) (bit \(b\) set).
Every \(x\in L, y\in R\) satisfies \(x\oplus y \ge 2^b\); every intra-group xor is
\(< 2^b\).  So the MST of \(S\) uses \(|L|-1\) edges inside \(L\), \(|R|-1\) inside
\(R\), and exactly one cross edge.
\end{observation}
\begin{observation}
Therefore \(\mathrm{MST}(S)=\mathrm{MST}(L)+\mathrm{MST}(R)+\min_{x\in L,y\in R}
(x\oplus y)\): a divide-and-conquer, and the cross term is the only non-recursive
work.
\end{observation}
\begin{observation}
\(\min_{x\in L,y\in R}(x\oplus y)\) is the minimum xor of one query value against
a fixed set -- answered greedily in \(\bigO{30}\) by inserting \(L\) into a binary
trie and walking, for each \(y\in R\), toward the same-bit child whenever possible.
\end{observation}

\solstep{Think It With Me}
\begin{thinkbox}
\deadend{Explicit Kruskal on all \(\binom n2\) edges}{\(\sim 2\cdot 10^{10}\)
edges at \(n=2\cdot10^5\): the sort alone is impossible.}
\deadend{Pair up equal values first and hope to shortcut}{Zero-xor pairs only help
duplicates; a graph with all-distinct values still needs the full structure, and
greedy local pairing ignores the cross-group cut.}
\aha{The highest differing bit forces the shape: recurse into the two bit-halves
(cut property pays one cross edge) and compute that one edge's xor with a trie.}
\end{thinkbox}

\solstep{The Approach Ladder}
\begin{center}\renewcommand{\arraystretch}{1.25}\setlength{\tabcolsep}{3pt}
\begin{tabular}{@{}c p{0.28\linewidth} p{0.30\linewidth} l c@{}}
\toprule
\textbf{\#} & \textbf{Approach} & \textbf{Idea} & \textbf{Time} & \textbf{Verdict}\\
\midrule
1 & \lv{3}\ Kruskal on all edges & materialise \(\binom n2\) xors & \(\bigO{n^2}\) & \mle\\
2 & \lv{4}\ greedy duplicate pair & shrink equal values & \(\bigO{n}\) & \textbf{WA}\\
3 & \lv{5}\ \textbf{bit D\&C + trie cross} & split on top bit, one cross & \(\bigO{30n}\) & \ac\\
\bottomrule
\end{tabular}\end{center}
Rung 2 is a tempting simplification that only helps when many values coincide; it
fails the second sample's structure.  Rung 3 is the real cut-property recursion.

\solstep{Solution}
A recursive function on a value multiset and a bit index.  Partition by the current
bit; if all values fall in one child, recurse into the other bit with no cross
cost.  Otherwise recurse on both halves and add the minimum cross xor, found by
inserting one half into a binary trie and querying every value of the other.

\solstep{Implementation}
\begin{lstlisting}[caption={CF 888G. Divide the value set on the highest bit; sum the two halves plus the minimum cross xor from a binary trie. Recursion is on bit depth (30), values move to exactly one child per level, so total work is O(30 n)}, label={lst:c17-summit}]
#include <bits/stdc++.h>
using namespace std;
typedef long long ll;
static const int HI = 29;                 // a_i < 2^30
int tr[6400005][2];                       // binary trie nodes ( <= 30 * n )
int trn;
void trieClear(){ trn=1; tr[0][0]=tr[0][1]=0; }
void trieInsert(int x){
  int cur=0;
  for(int b=HI;b>=0;b--){
    int d=(x>>b)&1;
    if(!tr[cur][d]){ tr[cur][d]=trn; tr[trn][0]=tr[trn][1]=0; trn++; }
    cur=tr[cur][d];
  }
}
int trieMinXor(int x){
  int cur=0, res=0;
  for(int b=HI;b>=0;b--){
    int d=(x>>b)&1;
    if(tr[cur][d]) cur=tr[cur][d];        // keep the bit, xor 0
    else { cur=tr[cur][d^1]; res |= (1<<b); }
  }
  return res;
}
ll solve(vector<int>& a, int bit){
  if(bit<0 || a.size()<=1) return 0;
  vector<int> L, R;
  for(int x : a){ ((x>>bit)&1 ? R : L).push_back(x); }
  if(L.empty()) return solve(R, bit-1);
  if(R.empty()) return solve(L, bit-1);
  trieClear();
  for(int x : L) trieInsert(x);
  int best=INT_MAX;
  for(int y : R) best=min(best, trieMinXor(y));   // min cross xor, full value
  return solve(L, bit-1) + solve(R, bit-1) + (ll)best;
}
int main(){
  int n; if(scanf("%d",&n)!=1) return 0;
  vector<int> a(n);
  for(int i=0;i<n;i++) scanf("%d",&a[i]);
  printf("%lld\n", solve(a, HI));
  return 0;
}
\end{lstlisting}

\solstep{Trace on the Official Samples}
\begin{tracebox}[Sample 1: \, the list \(1,2,3,4,5\)]
Top bit \(2\) (\(2^2=4\)): \(L=\{1,2,3\}\), \(R=\{4,5\}\).
\(\mathrm{solve}(R)\): bit0 splits \(\{4\},\{5\}\), cross \(4\oplus5=1\) \(\to 1\).
\(\mathrm{solve}(L)\): bit1 splits \(\{1\},\{2,3\}\); \(\{2,3\}\to\) cross
\(2\oplus3=1\); cross between \(\{1\}\) and \(\{2,3\}\) is \(\min(1\oplus2,1\oplus3)
=\min(3,2)=2\); total \(0+1+2=3\).  Cross \(L\)-\(R\) at bit2: \(1\oplus5=4\) is the
minimum.  Grand total \(3+1+4=\mathbf{8}\).
\end{tracebox}
\begin{tracebox}[Sample 2: \, the list \(1,2,3,4\)]
Bit2: \(L=\{1,2,3\}\) (solve \(=3\) as above), \(R=\{4\}\) (solve \(=0\)).  Cross
\(L\)-\(R\): \(\min(1\oplus4,2\oplus4,3\oplus4)=\min(5,6,7)=5\).  Grand total
\(3+0+5=\mathbf{8}\).
\end{tracebox}

\solstep{Complexity}
\begin{complexitybox}
Each value descends the bit-recursion exactly once per level, so the partitioning
across all nodes is \(\bigO{30n}\).  Every trie insert and query is \(\bigO{30}\),
and the cross term inserts one half and queries the other at each split --
amortised \(\bigO{30n}\) total.  Overall \(\bigO{30n}\) time and \(\bigO{30n}\)
trie nodes plus the recursion's temporary vectors: for \(n=2\cdot10^5\) that is
a few million operations, inside 2 seconds.  Sums fit in \texttt{long long};
individual xors are \(<2^{30}\).
\end{complexitybox}

\begin{hardeningbox}
\begin{itemize}[nosep,leftmargin=1.4em]
\item \textbf{Recursion depth.}  The recursion depth is the bit count (\(\le 30\)),
      not \(n\) -- safe on the system stack.  Do not, however, recurse on vector
      copies carelessly: total vector traffic is \(\bigO{30n}\); passing by value at
      every level is the constant factor to watch.
\item \textbf{Trie size.}  At most \(30n+1\) nodes; allocate the array globally (as
      here) rather than per call to avoid reallocation and stack pressure.
\item \textbf{Base cases.}  \(a.\mathrm{size()}<=1\) or \(\text{bit}<0\) returns
      \(0\); a one-element half contributes no internal edge, and the cross term is
      skipped whenever one side is empty.
\item \textbf{Memory of xors.}  \(\mathit{best}\) starts at \texttt{INT\_MAX} only
      because a cross edge always exists when both halves are non-empty; do not
      reuse the routine on an empty side.
\end{itemize}
\end{hardeningbox}

\begin{recallbox}[Chapter 31 in ten seconds]
\emph{One lemma rules the chapter: the lightest edge across any cut is safe
(\(\thmref{c17-cut}\)).  Kruskal and Prim are two machines for honouring it; the
reconstruction tree replays it so a bottleneck query is an LCA weight
(\(\thmref{c17-krtmax}\)); the classification asks which edges the cut ever allows
(\(\thmref{c17-intersect}\), \(\thmref{c17-cycle}\)); the second-best tree is the
cut property one swap away (\(\thmref{c17-second}\)); and the Summit runs the same
cut over value bits with a trie.  CF 1095F is a reduced Kruskal, CF 1108F is
Kruskal grouped by weight, CF 888G is Kruskal's idea without the edges.}
\selftest{
  \item State the cut property and name the two build algorithms it justifies.
  \item Why does a bottleneck (minimax) query become an LCA weight in the KRT?
  \item Give the per-equal-weight-group formula behind CF 1108F's answer.
  \item Sketch the Summit recursion: what does the cut property contribute at each bit split?
}
\end{recallbox}

\chaptersection{Tool Kit Inventory}

\begin{center}
\renewcommand{\arraystretch}{1.2}\setlength{\tabcolsep}{3pt}
\begin{tabular}{@{}p{0.26\linewidth} >{\raggedright\arraybackslash}p{0.11\linewidth} c >{\raggedright\arraybackslash}p{0.36\linewidth}@{}}
\toprule
\textbf{Tool} & \textbf{Since} & \textbf{Lvl} & \textbf{Trigger phrase}\\
\midrule
Cut property & Ch.31 & 3 & "lightest edge across a cut is safe"\\
Kruskal + DSU & Ch.31 & 4 & sparse graph, cheapest edges connect all\\
Prim + heap / matrix & Ch.31 & 4 & dense graph / adjacency matrix\\
Reconstruction tree & Ch.31 & 4 & "minimax / bottleneck path"\\
Edge classification & Ch.31 & 5 & "in some / every / no MST"\\
Group excess count & Ch.31 & 5 & "make the MST unique"\\
Second-best swap & Ch.31 & 4 & "cheapest tree besides the MST"\\
Trie bit D\&C & Ch.31 & 5 & "complete graph, xor weight"\\
DSU (black box) & Ch.17 & 3 & connectivity, Kruskal's union\\
LCA (black box) & Ch.16 & 4 & bottleneck value on KRT path\\
Bottom-up fold & Ch.13 & 4 & Summit's divide-and-conquer\\
Matrix-Tree count & Ch.30 & 4 & hand-off: "how many", not "cheapest"\\
Bridges / articulation & Ch.29 & 4 & hand-off: unweighted robustness\\
\bottomrule
\end{tabular}
\end{center}

\chaptersection{Related Topics}

\begin{itemize}[nosep]
  \item Ch.~17 (\(\chref{c28}\)) DSU -- Kruskal's \texttt{find}/\texttt{union}
        engine, used here as a black box and never re-derived.
  \item Ch.~16 (\(\chref{c11}\)) LCA -- the query layer for the Kruskal-reconstruction
        tree's bottleneck answers (\(\secref{c17-krt}\)).
  \item Ch.~13 (\(\chref{c10}\)) tree DP -- the bottom-up divide-and-conquer fold the
        Summit reuses over value bits.
  \item Ch.~30 (\(\chref{c32}\)) counting trees -- the enumeration twin: it counts
        spanning trees, this chapter chooses the cheapest.
  \item Ch.~29 (\(\chref{c31}\)) trees from graphs -- the unweighted robustness side
        (bridges, near-trees) that a "spanning" lookalike often really wants.
  \item Ch.~19 (\(\chref{c13}\)) segment trees -- cited from the earlier arc; not
        required inside Chapter 31.
  \item Ch.~12 (\(\chref{c09}\)) tree metrics -- distances/diameters on the resulting
        tree, e.g.\ the diameter of an MST.
\end{itemize}

\begin{lastword}[Last Word --- Chapter 31]
Chapter~31 reduces a whole family of "cheapest way to connect everything" problems
to a single sentence: the lightest edge crossing any cut is safe.  That sentence is
deceptively cheap -- it takes one paragraph to state and one lemma to prove, yet
Kruskal, Prim, the reconstruction tree, the edge trichotomy, the second-best swap,
and even the xor Summit are the same sentence seen from six angles.  The craft is
never in re-proving the lemma; it is in recognising which cut a problem is secretly
asking you to respect.  Where Chapter~30 taught you to \emph{count} the trees, this
chapter teaches you to \emph{choose} one -- and, once chosen, to interrogate its
paths, its ties, and its alternatives as fluently as you read a code off a tree.
Train three reflexes: see a weighted connect-all task and reach for Kruskal or Prim
by density; see a "heaviest edge on the best path" and build the reconstruction tree;
see "is this edge in an MST / is the MST unique" and split into equal-weight groups.
Two habits outlast any algorithm here: reason about every greedy from a cut, not from
a hunch; and never re-derive DSU or LCA when the book has already banked them
(\(\chref{c28}\), \(\chref{c11}\)) -- compose, don't rebuild.  You have now met the
tree from every direction the book set out to cover: shape, traversal, DP,
decomposition, isomorphism, extraction, enumeration, and finally optimisation.  The
next time a judge hands you a weighted graph and asks for the cheapest network,
you will already know the one line that decides it.  \begin{center}\(\blacksquare\)\end{center}
\end{lastword}
